\chapter{Literature Review}
\label{ch:literature-review}

\section{Preliminaries}

Chapter~\ref{ch:introduction} laid out the three limitations of current LLMs that CAEM addresses (stateless operation, poor confidence calibration, absence of self-improvement) and the three architectural commitments (verified episodic memory, confidence-aware routing, constrained self-improvement) that answer them. This chapter supplies the two kinds of background a reader needs to evaluate those commitments in detail. The Preliminaries section that follows establishes foundational concepts and terminology: the transformer backbone, the hallucination taxonomy, natural language inference, sentence embeddings, self-consistency decoding, and catastrophic forgetting. The Review of Existing Research section that follows it then surveys the empirical literature on hallucination measurement, verification, confidence estimation, memory-augmented architectures, selective retrieval, self-improvement and continual learning, external baselines, and evaluation benchmarks, locating CAEM's contribution relative to prior work. Together they prepare the ground for the formal methodology in Chapter~\ref{ch:methodology}.

\subsection{Large Language Models and Transformer Architecture}

Large language models are neural networks trained on very large text collections to predict and generate natural language. Modern large language models adopt the Transformer architecture \cite{NIPS2017_3f5ee243}, whose self-attention mechanism lets every token in the input attend to every other token. The original Transformer stacks encoder and decoder layers, each combining multi-head self-attention with a feed-forward network; the decoder-only variant used in this thesis keeps only the decoder stack.

This thesis uses Qwen-2.5-3B-Instruct \cite{qwen25}, a three-billion-parameter decoder-only transformer that has already been instruction-tuned and preference-tuned before it enters the CAEM pipeline. Decoder-only models generate the answer token by token from left to right, with the question and the generated rationale sharing a single causal attention stream. The backbone is chosen for two reasons: it fits the whole self-improvement loop, including low-rank fine-tuning, Monte Carlo dropout sampling and the nine-signal verifier, on a single consumer-grade graphics processor in sixteen-bit brain-float precision; and because it is instruction-tuned, it already answers in the required format, so the verifier has usable outputs to score from the first cycle.

Fine-tuning adapts pre-trained models to specific tasks or domains by continuing training on task-specific data with a smaller learning rate. The model parameters $\boldsymbol{\theta}$ are adjusted to reduce loss on the new data while preserving the general abilities learned during pre-training. Standard fine-tuning updates every parameter; parameter-efficient approaches keep the underlying backbone frozen and adapt only a small fraction of the parameter count through trainable side-modules. CAEM uses the parameter-efficient path through the low-rank adaptation (LoRA) layer of Hu et al.~\cite{hu2022lora} on the attention and feed-forward projections, with the trainable footprint at roughly two percent of the backbone parameter count and the underlying backbone frozen across every cycle of the self-improvement loop. The empirical evidence motivating that choice over full-parameter fine-tuning is reviewed in \Cref{sec:lit-self-improvement}.

\subsection{Hallucination Taxonomy}
\label{sec:hallucination-taxonomy}

Hallucinations in language model outputs can be understood along two main dimensions. The first dimension separates \textit{intrinsic hallucinations}, which directly contradict information provided in the source input, from \textit{extrinsic hallucinations}, which introduce information that can be neither supported nor contradicted by the source \cite{ji2023survey}. For instance, if a biography states that ``Einstein was born in Germany'' and the model generates ``Einstein was born in France,'' this is an intrinsic hallucination because it contradicts the stated source. But if the model generates ``Einstein played violin exceptionally well'' without any evidence in the source, this is an extrinsic hallucination even if the claim happens to be true: the defining characteristic is that it cannot be verified from the input, not that it contradicts anything in it.

The second dimension distinguishes \textit{factuality hallucinations}, which conflict with verifiable world knowledge, from \textit{faithfulness hallucinations}, which contradict user input or task context \cite{huang2025survey}. For example, a model that states ``photosynthesis produces oxygen from carbon dioxide and water'' is factually accurate, whereas one that states ``photosynthesis occurs in mitochondria'' produces a factuality hallucination. Faithfulness focuses on whether the output of the model remains consistent with the supplied input or context, especially when summarising documents or following instructions.

This research focuses on factuality hallucinations, which occur in knowledge-based reasoning tasks. Here, correct answers can be verified against trusted sources or logical rules. Faithfulness to retrieved evidence also matters for retrieval-augmented generation, where outputs must be grounded in retrieved documents.

Within factuality hallucinations, four subclasses are distinguished operationally because the four signal families of CAEM's verifier map onto them.\footnote{The nine signals are defined formally in \Cref{ch:methodology}; this chapter names them only where needed for motivation.} \textit{Confabulations}, in the sense of Farquhar et al.~\cite{farquhar2024semantic}, are answers the model might equally have produced differently on a resample, even when the individual token probabilities look confident; they are targeted by the sample-set agreement signals, which compare several resampled reasoning chains with one another ($s_{\text{avg}}$) and with the served answer ($p_{\text{entail}}$). \textit{Ungrounded factual claims} are answers without support in retrieved evidence; they are targeted by the three external-grounding signals, which score the answer against retrieved passages. \textit{Off-topic answers} that drift from the question are targeted by the question-answer relevance signal $q_{\text{a,rel}}$. \textit{Confidently wrong} answers, internally confident but externally unsupported, are caught by the confabulation early-exit gate, which fires when a high internal-confidence signal coincides with low grounding. Outside this scope, CAEM does \textit{not} target instruction-following failures, intermediate-reasoning-step errors inside chain-of-thought traces (only the final answer is verified), intrinsic faithfulness hallucinations in summarisation or translation (where the grounding source is the input document rather than world knowledge), or input-contradiction hallucinations that conflict with user-supplied premises. These exclusions set the failure-mode boundary within which the results in \Cref{ch:evaluation} should be interpreted.

\subsection{Natural Language Inference}

Natural Language Inference (NLI) is the task of determining whether a hypothesis sentence logically follows from a premise sentence \cite{barhaim2006second}. The PASCAL Recognising Textual Entailment challenges framed this as a binary decision (entailed or not entailed); later NLI corpora, including the Multi-Genre NLI corpus on which RoBERTa-Large-MNLI is fine-tuned, use three classes. Given premise $P$ and hypothesis $H$, a three-class NLI model classifies the relationship as:
\begin{itemize}
    \item \textbf{ENTAILMENT}: $P$ implies $H$ is true
    \item \textbf{CONTRADICTION}: $P$ implies $H$ is false  
    \item \textbf{NEUTRAL}: $P$ neither confirms nor denies $H$
\end{itemize}

For example, if the premise is ``The Beatles were formed in Liverpool in 1960'', the hypothesis ``The Beatles were an English band'' is entailed, the hypothesis ``The Beatles were formed in London'' is a contradiction, and the hypothesis ``The Beatles recorded many albums'' is neutral.

CAEM's verifier uses MiniCheck-Flan-T5-Large \cite{tang2024minicheck} as its entailment judge, a Flan-T5-Large checkpoint fine-tuned on synthetic claim-support data. MiniCheck is a binary judge: it returns the probability that a claim is supported by the given evidence and has no separate contradiction class, so a contradicted claim simply receives a low support probability. A second judge for long hypotheses was considered and dropped before the main run (\Cref{ch:methodology}), so every hypothesis is scored by MiniCheck, truncated at its 512-token input limit. Generic three-class NLI models such as RoBERTa-Large-MNLI \cite{liu2019roberta}, which reaches about 90.2\% on the MNLI matched development set, are the conventional alternative, but they transfer poorly to the hallucinations of language models: Sinha~\cite{sinha2025semantic} reports a one-hundred-percent false-positive rate at ninety-five-percent recall for a near-state-of-the-art generic NLI model on HaluEval~\cite{li-etal-2023-halueval}. Entailment checked in both directions can also decide whether two answers mean the same thing: Kuhn et al.~\cite{kuhn2023semantic} treat two answers as equivalent if $P$ entails $H$ and $H$ entails $P$. This is the basis of semantic entropy, which CAEM uses as an external baseline (\Cref{sec:external-baselines}).

\subsection{Sentence Embeddings and Semantic Similarity}

Sentence embeddings map variable-length text sequences to fixed-dimensional dense vectors in a semantic space where similar meanings have high cosine similarity. Sentence BERT (SBERT) \cite{reimers-gurevych-2019-sentence} produces sentence embeddings by fine-tuning BERT in a siamese network structure, first on natural language inference data (SNLI and MultiNLI) and optionally on semantic textual similarity data. The model used here, all-mpnet-base-v2 from the same Sentence-Transformers family (built on an MPNet backbone rather than the original BERT), generates 768-dimensional embeddings.

Cosine similarity between vectors $\mathbf{v}_1$ and $\mathbf{v}_2$ is computed as:
\begin{equation}
\text{sim}(\mathbf{v}_1, \mathbf{v}_2) = \frac{\mathbf{v}_1 \cdot \mathbf{v}_2}{\|\mathbf{v}_1\| \|\mathbf{v}_2\|} = \cos(\theta)
\end{equation}
where $\theta$ is the angle between vectors. Similarity ranges from $-1$ (opposite) to $1$ (identical), with $0$ indicating orthogonality.

Semantic similarity is different from string similarity. For instance, the questions ``Is 17 prime?'' and ``Is the number 17 a prime number?'' share few words but are close in meaning. CAEM uses semantic similarity for memory retrieval (finding similar stored questions), novelty detection (avoiding redundant storage), and the self-consistency signal (measuring agreement between reasoning chains). Nearest-neighbour search uses the FAISS library \cite{8733051}; at the memory sizes reached in this thesis the search is exact, and FAISS's approximate indexes take over only at much larger sizes.

\subsection{Self-Consistency Decoding}
\label{sec:lit-self-consistency}

Self-consistency \cite{wang2023selfconsistency} improves reasoning by generating several independent reasoning chains for the same question and selecting the most common answer. Unlike greedy decoding, which produces a single output, self-consistency samples $M$ chains at a temperature $T > 0$ to introduce diversity.

CAEM adapts self-consistency for verification rather than answer selection. Given a question $q$ and a generated answer $a$, the verifier generates $M$ additional chains and measures their agreement as the average pairwise similarity
\begin{equation}
s_{\text{avg}} = \frac{2}{M(M-1)} \sum_{1 \le i < j \le M} \text{sim}(\mathbf{v}_i, \mathbf{v}_j)
\end{equation}
where $\mathbf{v}_i$ is the embedding of chain $i$ and only distinct pairs are counted. High agreement suggests that several reasoning paths converge on the same content; low agreement suggests uncertainty. CAEM pairs $s_{\text{avg}}$ with the chain-to-answer entailment signal $p_{\text{entail}}$, computed on the same $M{=}3$ chains so that the decoding cost is paid once; together they form the sample-set agreement family of the verifier. The normalised semantic entropy of Farquhar et al.~\cite{farquhar2024semantic}, computed over $K{=}10$ samples at temperature $T{=}1.0$, was originally a third member of this family. It was dropped before the main run because its learned weight in the cycle-zero composite fit was effectively zero (about $-0.0005$), so it added nothing beyond the other signals, while its ten extra samples per answer were the largest single cost in the verifier (about four seconds per answer). The kernel-language-entropy generalisation of Nikitin et al.~\cite{nikitin2024kle}, a candidate replacement that needs the same extra samples, was not adopted either.

\subsection{Catastrophic Forgetting and Continual Learning}

Catastrophic forgetting, first described by McCloskey and Cohen~\cite{MCCLOSKEY1989109}, occurs when neural networks learn new tasks and suddenly lose their ability to perform previously learned tasks. Standard gradient descent focuses only on improving performance for the current task and does not protect earlier knowledge, so parameter updates that help the new task tend to weaken the old ones. In large language models, fine-tuning on a new domain likewise erodes performance outside that domain, and Biderman et al.~\cite{biderman2024lora} found that the erosion grows with training duration.

Elastic Weight Consolidation (EWC) \cite{doi:10.1073/pnas.1611835114} addresses this by estimating, through the Fisher information, which parameters matter for earlier tasks, and adding a quadratic penalty that discourages large changes to them:
\begin{equation}
\mathcal{L} = \mathcal{L}_{\text{new}} + \frac{\lambda}{2} \sum_{i} F_i (\theta_i - \theta_i^*)^2
\end{equation}
where $F_i$ measures parameter importance, $\theta^*$ are the previous parameters, and $\lambda$ controls the strength of the penalty. A cheaper variant drops the Fisher weights and penalises all parameters equally, $\frac{\lambda}{2}\lVert \boldsymbol{\theta} - \boldsymbol{\theta}_{\text{prev}} \rVert^2$; Kirkpatrick et al.~\cite{doi:10.1073/pnas.1611835114} show that this uniform penalty is weaker than EWC because it protects all weights equally and leaves little capacity for the new task. CAEM uses neither penalty on its shipped path. It bounds drift between cycles with three other mechanisms instead: a strict quality filter on the training pool; low-rank adaptation \cite{hu2022lora}, which trains only about two percent of the parameters on top of a frozen backbone and forgets less than full fine-tuning (\Cref{sec:lit-self-improvement}); and a retention guard that rolls the adapter back, while keeping the episodic memory, whenever accuracy on held-out probe sets falls below a fixed fraction of its starting value. The details are given in \Cref{ch:methodology}.

The conceptual basis for understanding CAEM's nine-signal post-generation verifier, pre-routing confidence, episodic-memory architecture, and constrained self-improvement loop is established by the groundwork above. The following sections describe how prior research develops and validates these individual components, preparing for their integration in Chapter~\ref{ch:methodology}.

\section{Review of Existing Research}

Where the Preliminaries section above defined the concepts and notation (hallucination taxonomy, natural language inference, sentence embeddings, self-consistency, catastrophic forgetting), this section surveys the prior empirical and methodological literature on hallucination measurement, verification, confidence estimation, memory-augmented architectures, selective retrieval, self-improvement and continual learning, the external baselines used in this thesis, and evaluation benchmarks. Each subsection identifies key contributions and limitations that motivate CAEM's architectural choices, so that a reader who has already absorbed the Preliminaries can move directly to the comparative positioning without re-reading definitional material.

\subsection{Hallucination Problem and Measurement}

Understanding the nature and characteristics of hallucination in LLMs requires methodologies like taxonomic frameworks and quantitative measurement. Early studies of hallucination focused on summarisation and data-to-text generation. Ji et al. \cite{ji2023survey} published the first major survey in ACM Computing Surveys. It distinguishes intrinsic hallucinations, which contradict the source material, from extrinsic hallucinations, which can be neither supported nor contradicted by it. Their work shows that hallucinations arise at several stages, including data collection that embeds factual errors, training procedures that reinforce misleading correlations, and inference methods that fail to account for uncertainty.

Huang et al. \cite{huang2025survey} extended this study specifically for large language models in ACM Transactions on Information Systems, distinguishing factuality hallucinations from faithfulness hallucinations. Their analysis traces the causes of hallucination across the model lifecycle to data, training and inference, and ties each family of mitigation methods to the cause it addresses. This suggests that effective mitigation should intervene at several stages instead of relying solely on detection applied after the model has already produced an answer, which supports CAEM's multi-component strategy combining memory, verification, and iterative learning.

Quantitative measurement studies of hallucination prevalence reveal concerning statistics across domains and model scales. Lin et al. \cite{lin-etal-2022-truthfulqa} introduced TruthfulQA at ACL 2022, a benchmark of 817 questions across 38 categories designed to elicit imitative falsehoods where models reproduce common human misconceptions found in training data. Their result demonstrates an inverse scaling effect: the best GPT-3 175B configuration achieved only 21\% on the truthful-and-informative metric (and 58\% on truthfulness alone), against 87\% and 94\% respectively for the human baseline. Larger models often performed worse than smaller ones. This suggests that scaling alone is not sufficient for truthfulness, which motivates architectural innovation.

The inverse-scaling result has a direct implication for mitigation: if adding parameters can reduce truthfulness on some tasks, mitigation must work beyond parameter growth. CAEM does so by adding external verification and memory components whose operation does not depend on the size of the model.

Min et al.\ \cite{min-etal-2023-factscore} introduced FActScore at EMNLP 2023, which evaluates long-form generations by decomposing the answer into atomic claims and reporting the share supported by a reliable knowledge source rather than scoring the answer as a whole. Its automated estimator stays within an error of 2\% of human evaluation, and on biography generation it finds that only 58\% of ChatGPT's atomic facts are supported, so even highly capable models produce many unsupported claims. CAEM borrows this decomposition for one of its grounding signals, $p_{\text{ground\_atomic}}$, which splits the answer into atomic claims, scores each against the retrieved passages, and keeps the weakest score. The step runs only on answers long enough to contain several claims; the short answers of the five-benchmark panel almost never qualify, so in the realised run this signal coincided with the passage-mean grounding signal on every evaluated answer (\Cref{ch:methodology}).

Domain-specific studies show even higher hallucination rates in specific contexts. Alkaissi and McFarlane \cite{alkaissi2023artificial} asked ChatGPT to explain a clinical finding and to supply supporting references; none of the five paper titles it returned existed. Chelli et al.\ \cite{info:doi/10.2196/53164} checked the references that chatbots generated for systematic reviews and reported hallucination rates of 39.6\% for GPT-3.5, 28.6\% for GPT-4, and 91.4\% for Google Bard.

Three points follow for CAEM's design. First, hallucination occurs across models, tasks and domains, at rates that vary widely with context (from 28.6\% to 91.4\% of generated references in the single study of Chelli et al.~\cite{info:doi/10.2196/53164}). Second, larger models do not automatically hallucinate less, so the structure around the model has to change. Third, scoring individual claims against evidence is more informative than judging whole responses, which motivates CAEM's evidence-based grounding signals.

\subsection{Verification Methods}

Verification methods are used to check whether a model's reasoning is correct before the information is stored or shown to users. This subsection reviews chain-of-thought prompting, self-consistency decoding, semantic entropy, natural-language inference and conformal filtering, the approaches from which CAEM's verifier draws its signals or against which its design is positioned.

Chain of thought prompting, which was introduced by Wei et al. \cite{wei2022chain} at NeurIPS 2022, improves reasoning quality by prompting models with worked examples whose intermediate steps precede the final answer. Showing these intermediate steps leads to stronger reasoning and provides better starting points for verification. When tested on the PaLM 540B model, chain of thought greatly improved performance on math reasoning tasks, increasing accuracy on GSM8K from 17.9\% to 56.9\%. The authors found that the technique is an emergent ability of model scale: below roughly 100 billion parameters chain-of-thought prompting did not help and could lower performance, a pattern that held for the instruction-tuned InstructGPT models in their GPT-3 series as well. Kojima et al.~\cite{kojima2022large} later showed that the single trigger ``Let's think step by step'' elicits chain-of-thought reasoning without worked examples. Instruction tuning does not remove the small-scale effect: Chung et al.\ \cite{JMLR:v25:23-0870} report that their instruction-tuned Flan-T5-Large (780 million parameters) still scores lower with chain-of-thought prompting than with direct answering on MMLU (40.5 versus 45.1 percent) and on BIG-Bench Hard (31.5 versus 37.5 percent). Whether chain-of-thought prompting helps CAEM's three-billion-parameter instruction-tuned backbone, Qwen-2.5-3B-Instruct~\cite{qwen25}, is therefore an empirical question, which the chain-of-thought baseline in \Cref{ch:evaluation} answers.

For CAEM, chain-of-thought prompting matters in two ways: it is the reasoning format in which the generator answers, and it supplies the resampled reasoning chains from which the sample-set agreement signals are computed. The verifier, however, scores only the final answer and not the individual reasoning steps (\Cref{sec:hallucination-taxonomy}), so an invalid step that still reaches a correct conclusion goes undetected.

Wang et al.\ \cite{wang2023selfconsistency} extended this idea in 2023 by introducing self-consistency. Instead of producing a single reasoning path, the model generates several different explanations and then selects the most common answer. The key idea is that difficult and complex problems usually have multiple valid reasoning paths that lead to the same correct answer. On the other hand, incorrect answers rarely achieve consistency across independent derivations. Reported gains over chain-of-thought prompting reach +17.9 percentage points on GSM8K, +11.0 on SVAMP, +12.2 on AQuA, and +6.4 on StrategyQA; these are the largest gains across the four models tested, and for PaLM 540B the corresponding gains are +17.9, +7.6, +12.5 and +6.3.

CAEM uses self-consistency differently: instead of selecting the majority answer, it measures how far the resampled chains agree in meaning. Strong agreement raises the composite score and weak agreement lowers it. The cost is a small number of extra generations per answer, which keeps the signal practical for routine verification.

Semantic entropy, introduced by Farquhar et al.\ \cite{farquhar2024semantic} in Nature in 2024, addresses a limitation of token-level entropy: answers that express the same idea in different words are grouped together by bidirectional entailment, and uncertainty is measured across these meaning-based groups rather than across surface strings.

Averaged over thirty combinations of tasks and models (LLaMA~2, Falcon and Mistral families), semantic entropy reaches an AUROC of $0.790$ for detecting confabulations, ahead of naive entropy ($0.691$), P(True) ($0.698$), and an embedding-regression baseline ($0.687$). The method targets arbitrary, sample-to-sample errors: a model that is consistently wrong gives the same answer on every sample and therefore looks certain.

CAEM originally included semantic entropy as a verifier signal and dropped it before the main run (\Cref{sec:lit-self-consistency}); it remains in the thesis as an external baseline that abstains when entropy is high (\Cref{sec:external-baselines}).

Natural language inference provides direct verification by checking whether a claim is supported by evidence. Thorne et al.\ \cite{thorne-etal-2018-fever} introduced the FEVER dataset at NAACL 2018, a fact-verification task closely related to NLI, with 185,445 claims generated from Wikipedia; unlike NLI, the evidence for each claim must first be retrieved rather than being given. Systems must classify each claim as SUPPORTS, REFUTES or NOT ENOUGH INFO.

Pre-trained NLI models such as RoBERTa-Large-MNLI reach about 90\% accuracy on MNLI and verify individual factual statements well, but they can miss reasoning errors that contradict no fact. A model could state ``Sydney is Australia's largest city'', which is true, and conclude ``therefore Sydney is the capital'', an invalid inference; NLI would confirm the premise and miss the faulty step, which is why complementary signals are needed.

A more recent strand of work brings conformal prediction to the language-model claim level. Mohri and Hashimoto~\cite{mohri2024conformal} introduced \emph{Language Models with Conformal Factuality Guarantees} at ICML 2024, decomposing each generated answer into atomic sub-claims, scoring each, and removing claims below a calibrated threshold $\tau$ such that the surviving sub-claims are jointly correct with probability at least $1-\alpha$ on a small annotated calibration set. On FActScore their method raises factuality from about 30\% to 80\% while retaining most sub-claims. Yadkori et al.~\cite{yadkori2024abstention} extended this with conformal-risk-control to bound the hallucination rate at a user-specified target through learned abstention rather than claim removal, providing the direct precedent for CAEM's abstention class in the four-outcome decision tree. Quach et al.~\cite{quach2024conformal} present \emph{Conformal Language Modeling} at ICLR 2024, calibrating a stopping rule, a rejection rule, and a confidence rule simultaneously through Learn-then-Test multiple-hypothesis testing. Cherian et al.~\cite{cherian2024conformal} sharpen the Mohri--Hashimoto framework with conditional guarantees and a boosted score, learning a linear combination of several per-claim scores by differentiating through the conformal procedure; the closest precedent for CAEM's boost layer over the per-signal log-odds in \Cref{sec:composite-boost} is the logistic-regression multi-scoring of Valentin et al.~\cite{valentin2024hallucination}. CAEM considered split-conformal calibration for its storage gate and rejected it because the distribution shift between its calibration and evaluation folds broke the coverage premise (\Cref{sec:conformal-gate}). The deployed gate is therefore a fixed-threshold rule on the calibrated probability, in the lineage of the post-hoc calibration analyses of Niculescu-Mizil and Caruana~\cite{niculescu2005predicting} and the selective-classification framing of Geifman and El-Yaniv~\cite{geifman2017selective}.

These methods supply complementary signals to CAEM's nine-signal verifier. Self-consistency and chain-to-answer entailment form the sample-set agreement family ($s_{\text{avg}}$, $p_{\text{entail}}$). Entailment against retrieved evidence forms the external-grounding family: the best-passage score $p_{\text{ground\_max}}$, the mean-passage score $p_{\text{ground\_mean}}$, and the length-gated atomic-fact score $p_{\text{ground\_atomic}}$ in the FActScore convention~\cite{min-etal-2023-factscore}. Three internal-calibration signals (token probability, Monte Carlo dropout variance, and an internal-state composite of the two) and the question-answer relevance signal complete the nine. Instead of a fixed weighted sum, each signal is mapped to a calibrated probability of correctness by isotonic regression~\cite{niculescu2005predicting}, and the calibrated signals are combined by a regularised logistic regression in the manner of Valentin et al.~\cite{valentin2024hallucination}; the full specification is in \Cref{ch:methodology}. The point for this chapter is that the literature supplies families of signals whose failure modes differ, which justifies a calibrated multi-signal composite over any single-signal detector.

\subsection{Confidence Estimation}

Reliable confidence estimation helps systems tell the difference between answers that can be trusted and ones that need extra checking from outside sources. This subsection looks at basic methods for measuring uncertainty and explains how these ideas have been adapted for use with large language models.

Gal and Ghahramani \cite{pmlr-v48-gal16} laid the theoretical foundation for practical uncertainty estimation in their ICML 2016 paper ``Dropout as a Bayesian Approximation''. They showed that dropout, normally used only during training, can be read as approximate Bayesian inference over the model's parameters. Keeping dropout active at test time and running the model several times then measures how much the outputs vary. This variation estimates epistemic uncertainty, which comes from the model's limited knowledge, as opposed to aleatoric uncertainty, which comes from noise in the data and cannot be reduced by more knowledge.

The framework shows that a network with dropout before each weight layer approximates a deep Gaussian process, whose predictive mean and variance can be estimated by Monte Carlo sampling: $K$ forward passes with different dropout masks, followed by the variance of the outputs. The estimate needs no ensemble training and no change to the loss function, and it applies to any network that already uses dropout, which makes it convenient for existing large language models.

For CAEM, MC Dropout is one of the three internal-calibration signals inside the post-generation verifier ($u_{\text{token}}$, $u_{\text{dropout}}$, and the derived $u_{\text{internal}} = 0.5\,u_{\text{token}} + 0.5\,(1 - u_{\text{dropout}})$ that combines them). It quantifies model uncertainty through output variance: high variance indicates the model is unstable with respect to its own weights and suggests that the composite score should be pulled down, while low variance indicates confident generation. This signal is insufficient on its own, since models can be confidently wrong, but the computational cost of the registered five forward passes is acceptable inside the verifier, which only runs on queries that have already been dispatched to Tier~2 or Tier~3 generation. MC~Dropout is deliberately \textit{not} used in the pre-routing confidence $\upre$, which must be cheap enough to run on every query before tier selection.

A major problem with neural network confidence is that models are often too confident in their answers. Guo et al. \cite{pmlr-v70-guo17a} demonstrated this at ICML 2017. They showed that modern deep networks predict confidence scores that substantially exceed their actual accuracy; a ResNet-110 on CIFAR-100, for example, has an expected calibration error of 16.5 percentage points before calibration. They popularised the Expected Calibration Error (ECE) of Naeini et al.~\cite{naeini2015obtaining} as the standard metric, dividing predictions into bins by confidence and measuring the gap between confidence and accuracy within each bin.

Their proposed solution, temperature scaling, achieves remarkable effectiveness with minimal complexity. The method divides logits by a learned scalar temperature $T > 0$ optimised on a validation set: $\hat{y} = \sigma(\mathbf{z}/T)$ where $\mathbf{z}$ are logits and $\sigma$ is softmax. Temperature scaling preserves accuracy (argmax unchanged) of classification while calibrating probabilities. Even though only a single parameter is used, it outperforms more complex methods including histogram binning, isotonic regression, and Bayesian approaches. The simplicity and effectiveness make temperature scaling the standard calibration method applied after training is complete.

CAEM applies temperature scaling to its pre-routing confidence $\upre$, fitted per benchmark on a held-out calibration split, so that a stated confidence of eighty percent corresponds to roughly eighty percent correct answers. The post-generation composite $\ustored$ is calibrated differently, by per-signal isotonic regression and a logistic combination (\Cref{ch:methodology}). Both calibrations make it possible to set meaningful thresholds for tier selection, for the four-outcome storage decision, and for the training-pool filter $\tau_{\text{train}}$.

Kadavath et al.\ \cite{kadavath2022language} from Anthropic explored self-evaluation in large language models in a 2022 preprint, introducing two ways of extracting confidence from a model itself. The first, P(True), asks the model to estimate the probability that its own earlier answer is correct. The second, P(IK), the probability that ``I know'' the answer, is predicted before any answer is proposed, by a model trained for that purpose. Their findings reveal nuanced calibration properties. Larger models tend to be well calibrated on multiple-choice questions, meaning that their confidence matches how often they are actually correct. However, policies fine-tuned with reinforcement learning from human feedback appear badly miscalibrated, because reinforcement learning collapses predictions toward the behaviours that receive the most reward, although a single temperature adjustment largely restores calibration.

The finding that reinforcement learning from human feedback can hurt calibration matters for deployment. A model trained to be helpful may express high confidence to sound authoritative even when it is uncertain, the opposite of what high-stakes use requires. CAEM therefore relies on several external confidence signals rather than on the model's own confidence alone.

Xiong et al.\ \cite{xiong2023can} studied confidence estimation for large language models in detail, comparing white-box methods, which use internal information such as token logits and sequence probabilities, with black-box methods, which rely on external signals such as verbalised confidence or consistency checks. Verbalised confidence is overconfident and discriminates poorly: even for GPT-4 its average AUROC for telling correct answers from incorrect ones is about 0.63, close to the 0.5 of random guessing. Sampling-consistency signals do much better on some tasks (up to 0.93 AUROC on GSM8K), but no single method consistently outperforms the others, and all of them struggle on tasks that require professional knowledge.

CAEM's design follows from this finding that no single signal is reliable across tasks, and splits confidence into two stages. The first, the pre-routing confidence $\upre$, fuses two cheap signals (the token probability of a short draft decode and a layer-convergence ratio over the query's hidden states) through a temperature-scaled sigmoid; it is computed for every query and drives tier selection and the safety override. The second is the nine-signal post-generation verifier, which runs only after zero-shot or retrieval-augmented generation and produces the calibrated composite $\ustored$ that drives the storage decision.

Recent work by Tian et al. \cite{tian2023just} asks how calibrated confidence can best be elicited from models fine-tuned with human feedback. Their study found that the raw token probabilities of such models are consistently overconfident, and that confidence verbalised as output tokens is typically better calibrated, often reducing the expected calibration error by a relative 50\%. Because the best confidence source differs between model families and tasks, this reinforces CAEM's strategy of combining implicit signals like token probability, dropout, and self-consistency with explicit ones, and of recalibrating every signal rather than trusting any raw score.

A complementary 2024 survey by Geng et al.~\cite{geng2024survey} catalogues white-box and black-box confidence-estimation methods for language models across the same axes that drive CAEM's verifier ensemble, and is the positioning document for CAEM's multi-signal design choice. Two recent single-signal methods deserve specific mention as alternatives considered for CAEM. Chen et al.~\cite{chen2024inside} introduced INSIDE / EigenScore at ICLR 2024, computing the log-determinant of the covariance matrix over middle-layer hidden states across $K$ resampled responses; the method reports AUROC of $0.80$ to $0.83$ on TriviaQA and clear gains over other unsupervised signals on CoQA, SQuAD and Natural Questions (on TriviaQA, perplexity scores slightly higher), and is registered in CAEM as a candidate replacement for the dropout signal that was not adopted in the locked configuration. Nikitin et al.~\cite{nikitin2024kle} introduced Kernel Language Entropy at NeurIPS 2024, generalising Farquhar et al.'s hard-cluster semantic entropy to a soft-kernel von Neumann entropy over pairwise semantic similarities; KLE was considered as a replacement should the semantic-entropy family return, but was not adopted because, like semantic entropy, it needs many extra samples per answer (\Cref{sec:lit-self-consistency}). The selective-prediction precision-versus-coverage tradeoff that CAEM reports in \Cref{ch:evaluation} follows the form pioneered by Kamath, Jia, and Liang~\cite{kamath2020selective} for selective question answering under domain shift, where a calibrator over feature signals predicts $P(\text{correct})$ and the selective rule answers only when the calibrator's confidence exceeds the registered precision target.

The confidence-estimation literature yields four points for CAEM. A single confidence signal is not reliable enough for safe routing, so several must be combined. Systematic overconfidence calls for calibration after training. White-box signals from inside the model complement black-box signals from output consistency. Preference tuning can degrade calibration, and CAEM's backbone is itself preference-tuned (its post-training uses DPO and GRPO \cite{qwen25}), so no raw confidence is used without recalibration. Together these points shape CAEM's two-stage confidence design: a temperature-scaled pre-routing score $\upre$ on every query, and a nine-signal verifier whose calibrated composite $\ustored$ drives the four-outcome storage decision. Both calibrations are fitted on labelled calibration data and refitted at every cycle boundary (\Cref{ch:methodology}).

\subsection{Memory-Augmented Architectures}

External memory architectures give neural networks a place outside their parameters where facts can be written and read. This subsection traces these ideas from early differentiable memories to retrieval-augmented generation.

Graves et al.\ \cite{graves2014neural} introduced Neural Turing Machines (NTMs) in 2014. NTMs gave a neural network an external memory that it can read from and write to in a differentiable, trainable way. The architecture combines a controller network with an addressable memory matrix, supporting both content-based addressing (via cosine similarity) and location-based addressing (via shift operations). The key innovation is a read/write mechanism that enables end-to-end training via backpropagation. As a result, NTMs successfully learned algorithmic tasks including copying, sorting, and associative recall.

For CAEM, the lesson is that an external memory can give a model abilities its parameters alone lack. NTMs were applied mainly to algorithmic tasks, but the same principle extends to facts: CAEM's episodic memory stores verified question-answer pairs together with their reasoning, so that a repeated question can be answered by lookup.

Weston et al. \cite{weston2015memory} introduced Memory Networks at ICLR 2015. Their model scaled to millions of memory slots. It also showed that a system could perform several rounds of reasoning to answer a question. The architecture consists of four learned components. The components are: turning the input into features, updating the memory, turning the memory output into features, and finally producing a response. Furthermore, Memory Networks performed strongly on bAbI question-answering tasks that required reasoning over multiple facts. Sukhbaatar et al. \cite{NIPS2015_8fb21ee7} extended this model with end-to-end Memory Networks at NeurIPS 2015. They removed the need for supervision on which memories to retrieve by using a continuous attention mechanism over the memory. This line of work motivates CAEM's use of stored knowledge as an explicit, inspectable memory.

Graves et al. \cite{graves2016differentiable} presented the Differentiable Neural Computer (DNC) in Nature 2016. It built on the idea of Neural Turing Machines but added smarter memory tools like dynamic allocation. It also keeps track of the order in which things were written. DNCs achieved a 3.8\% mean error on the bAbI tasks, compared with 7.5\% for the best previous jointly trained model, which suggests that learned memory management improves performance.

Lewis et al. \cite{NEURIPS2020_6b493230} introduced Retrieval-Augmented Generation (RAG) at NeurIPS 2020. They showed that combining what a model knows in its parameters with information stored outside the model can make its answers much more accurate. RAG adds a large searchable index of Wikipedia to a BART model with four hundred million parameters. It uses Dense Passage Retrieval to look up relevant documents. Then it works out the final answer by considering the generation probability across all the documents it retrieved. The architecture achieves substantial improvements over parametric-only models: 44.5\% exact match on Natural Questions (versus 41.5\% prior SOTA), 68.0\% on TriviaQA (versus 60.5\%), and 45.5\% on WebQuestions (versus 44.7\%).

Critically, RAG demonstrates index hot-swapping: updating the document index without retraining achieves 70\% accuracy on questions about updated world knowledge, proving that non-parametric memory can be updated without expensive retraining. This property directly motivates CAEM's episodic memory, which can be updated continuously through verification without model fine-tuning, though CAEM additionally employs iterative fine-tuning to consolidate learned patterns.

The limitation of standard RAG for CAEM's purposes is that retrieved knowledge is not verified before use and does not accumulate verified reasoning chains. CAEM extends RAG by adding verification layers and storing verified outputs for future retrieval, creating a growing knowledge base rather than relying solely on static external corpora.

Johnson et al.\ \cite{8733051} developed FAISS in 2019 for similarity search over billions of items, enabling efficient memory-augmented systems. The library achieves an 8.5$\times$ speedup over prior GPU methods through optimised k-selection, and supports index types such as an inverted file with product quantisation, which searches in sub-linear time. For CAEM, FAISS provides fast similarity search over the stored episodes, of which there were $12\,175$ at the end of the realised trajectory.

Reimers and Gurevych \cite{reimers-gurevych-2019-sentence} introduced Sentence-BERT (SBERT) at EMNLP 2019. SBERT produces semantically meaningful sentence embeddings through siamese-network fine-tuning. The speed improvement over standard BERT is substantial: finding the most similar sentence pair in a group of ten thousand sentences takes about sixty-five hours with standard BERT, but only about five seconds with SBERT, because one can compare the precomputed embeddings directly instead of running a heavy cross-encoder over every pair. This efficiency gain is what makes CAEM's memory search possible: each query is turned into a 768-dimensional embedding, and the system can look through the entire memory index in real time.

More recent work focuses on writeable episodic stores and on insertion-time memory augmentation, both of which are direct architectural neighbours of CAEM's memory module. Das et al.~\cite{larimar2024das} introduced Larimar at ICML 2024, an associative key-value episodic memory with one-shot edits, selective forgetting, and an order-of-magnitude faster knowledge-editing path than direct fine-tuning; Larimar establishes that a writeable episodic store is a feasible alternative to parametric storage but does not associate a quality scalar with each entry. Xu et al.~\cite{xu2025amem} introduced A-MEM, a Zettelkasten-inspired agentic memory that writes structured textual notes at insertion and evolves links between memories over time; A-MEM is the closest architectural neighbour to CAEM in the sense that both write structured metadata at insertion, but A-MEM's metadata is textual (tags, links) while CAEM's is a single calibrated scalar storage score that drives a tiered router rather than evolving link structure. Adaptive-RAG \cite{jeong2024adaptive}, reviewed in \Cref{sec:lit-selective-retrieval}, shares CAEM's count of three tiers but routes on query complexity at inference time rather than on the confidence stored with each memory entry, which is the distinction CAEM's router draws (\Cref{sec:router}).

These developments establish the feasibility and benefit of external knowledge storage. NTMs show that differentiable memory supports structured reasoning; Memory Networks scale to millions of facts with multi-hop reasoning; RAG shows that non-parametric knowledge improves factual accuracy and can be updated; FAISS and Sentence-BERT provide the infrastructure for practical deployment. CAEM builds on these advances in three ways. First, a query is routed by a combined score $\mathcal{S} = \lambda\, s + (1-\lambda)\, \ustored$ of its similarity $s$ to the nearest stored entry and that entry's stored confidence, so a close match stored with low confidence is not served from memory. Second, an independent safety override sends any query whose pre-routing confidence $\upre$ falls below a fixed floor to the retrieval-augmented default, however close its memory match, after which the retrieval-utility classifier makes the final choice between retrieval and zero-shot generation. Third, the memory is filled only through the four-outcome storage decision (\textsc{Store}, \textsc{Defer}, \textsc{Abstain}, \textsc{Discard}), with a bounded buffer for borderline answers, so the knowledge base grows only from answers that clear the storage threshold. These extensions distinguish CAEM from static-corpus RAG.

\subsection{Selective and Adaptive Retrieval}
\label{sec:lit-selective-retrieval}

The retrieval-augmented architectures of the previous subsection retrieve unconditionally on every query: every incoming question pays the retrieval cost, and the retrieved passages enter the generator's context whether the query benefits from external evidence or not. A complementary line of work argues that retrieval is not uniformly helpful and proposes architectures that decide whether to retrieve on a per-query basis. The selective-retrieval literature anchors the architectural design of \Cref{sec:ruc}, which adds a binary retrieval-utility classifier between CAEM's three-tier router and its retrieval-augmented fallback so that queries on which retrieval would amplify the question's misconception framing route to the zero-shot tier instead.

Mallen et al.~\cite{mallen2023whennot} established the empirical anchor for the selective-retrieval position at ACL 2023. They introduced PopQA, a fourteen-thousand-question benchmark of entity-centric factoid queries, and showed that large language models handle popular entities through their parametric memory more reliably than through retrieval, while the same models fail systematically on long-tail entities that retrieval can rescue. The architectural lesson is that retrieval utility is a function of the subject entity's popularity rather than a global hyperparameter, and the corresponding routing rule is a per-relation popularity threshold derived from Wikipedia pageview ranks. With GPT-3 davinci-003, the adaptive retrieval rule reaches 46.5\% accuracy on PopQA, 5.3 percentage points above the best non-adaptive method. The lesson for CAEM is that retrieval utility can be predicted in part from cheap question-side features such as entity popularity, which CAEM's retrieval-utility classifier includes among its inputs (\Cref{sec:ruc}).

Wang et al.~\cite{wang2023skr} extended this line to a learned classifier at EMNLP Findings 2023. Their self-knowledge routing system collects training labels by running each candidate question through both a no-retrieval pipeline and a retrieval-augmented pipeline, recording the per-pipeline exact-match score, and labelling the question as \emph{retrieval-helpful} when the retrieval-augmented exact match exceeds the no-retrieval exact match and as \emph{retrieval-hurt} otherwise. A classifier is then built on the labelled set (the paper compares direct prompting, in-context learning, a BERT-base classifier and a k-nearest-neighbour rule, of which the last performs best), and at inference the classifier dispatches each query to retrieval or to the zero-shot generator according to its prediction. The training-label rule (compare exact match across the two routings) is the direct precedent for the empirical-delta labelling rule that CAEM's retrieval-utility classifier inherits in \Cref{sec:ruc}; the architectural difference is that CAEM's classifier sits downstream of the existing tier-1 and tier-2 dispatch outcomes and gates only the retrieval-augmented fallback, while Wang's classifier replaces the dispatch entirely.

Jeong et al.~\cite{jeong2024adaptive} introduced Adaptive-RAG at NAACL 2024 as a three-class extension of the binary routing decision. Their classifier (a T5-large language model fine-tuned on automatically-collected complexity labels) dispatches each query to one of three classes (no retrieval, single-step retrieval, multi-step retrieval), trading routing accuracy for the cost of running a large model at routing time. The reported cost-accuracy frontier on a six-dataset open-domain question-answering suite (SQuAD, Natural Questions, TriviaQA, HotpotQA, 2WikiMultiHopQA, MuSiQue) beats single-strategy baselines and the unconditional retrieval-augmented baseline; for CAEM's purposes the T5-large routing cost is the design's main weakness, while the authors themselves single out the quality of the automatically assigned complexity labels. The architectural lesson for CAEM is that a learned classifier on automatically-collected outcome labels can decide the routing question, but that putting a large model at routing time costs more than the architecture gains; the follow-up of Marina et al.~\cite{moskvoretskii2025llm} extends Adaptive-RAG to the regime where the classifier carries no language-model forward pass and matches the routing accuracy of the language-model-based variant on the same suite. CAEM's retrieval-utility classifier sits at the lightweight end of this design space: the classifier itself is a logistic regression head, and its features reuse quantities the pipeline already computes, including a self-knowledge probe read from one forward pass of the base generator.

Asai et al.~\cite{asai2024selfrag} introduced Self-RAG at ICLR 2024 as the end-to-end fine-tuned alternative to the external-classifier approach. Self-RAG trains a single language model to emit four kinds of reflection tokens during generation (a \texttt{Retrieve} token that triggers retrieval, an \texttt{IsREL} token that judges passage relevance, an \texttt{IsSUP} token that judges support, and an \texttt{IsUSE} token that judges usefulness), so the routing decision is made by the same model that generates the answer. The deployed signal for retrieval routing is the probability the model assigns to the \texttt{Retrieve} token at decoding time; removing retrieval at inference time drops PopQA accuracy by roughly forty-five percent relative in the authors' ablation, the empirical receipt for the joint design's load-bearing role. The architectural cost is that the reflection tokens are tied to the LLaMA-family backbone Self-RAG fine-tunes on, and retargeting them to a different backbone changes enough of the system to invalidate the architectural comparison; this is the same constraint that prevents Self-RAG from entering the external baseline panel of \Cref{ch:evaluation} as a numerical baseline rather than a citation-level reference.

Baek et al.~\cite{baek2025probing} introduced Probing-RAG at NAACL Findings 2025 as a midpoint between the heavy classifier and the in-generator reflection tokens. They train a linear probe on the intermediate hidden states of the generator (rather than on its output tokens), with the probe trained under a balanced positive-negative protocol that explicitly enforces a fifty-fifty class ratio. The probe is trained to predict, from the hidden states, whether the model's current answer is correct and hence whether additional retrieval is needed; the reported accuracy across five open-domain question-answering datasets beats single-step retrieval, FLARE, DRAGIN, and Adaptive-RAG on the same suite while retrieving roughly half as often. The architectural lesson for CAEM is the most directly relevant of the five: the closest cousin of CAEM's retrieval-utility classifier on the feature axis is the hidden-state-as-routing-signal pattern that Probing-RAG establishes, with the methodological difference that Probing-RAG probes intermediate layers, whereas CAEM reads the base generator's last-token hidden state, without the adapter, as a single self-knowledge feature alongside thirty-five other question-side and passage-side features.

CAEM's retrieval-utility classifier reaches a pooled out-of-fold AUROC of $0.850$ under nested cross-validation on $3\,208$ labelled rows (\Cref{sec:ruc}). The systems reviewed above are evaluated on different question suites and with different metrics, so this figure is not directly comparable with theirs.

\subsection{Self-Improvement and Continual Learning}
\label{sec:lit-self-improvement}

Self-improvement mechanisms enable models to learn from interaction and feedback, while continual learning addresses the challenge of acquiring new knowledge without forgetting previous capabilities. This subsection reviews methods for iterative refinement and catastrophic forgetting mitigation that inform CAEM's self-improvement loop.

Kirkpatrick et al.\ \cite{doi:10.1073/pnas.1611835114} introduced Elastic Weight Consolidation (EWC), defined in the Preliminaries, in PNAS in 2017. Its Fisher-weighted penalty lets parameters that matter little for earlier tasks adapt freely while protecting those that matter. On ten Atari games learned in sequence, EWC let a single agent learn many games without catastrophic forgetting, whereas plain stochastic gradient descent forgot earlier games, although EWC did not reach the score of ten separately trained agents. For CAEM, EWC supplies the motivation for bounding drift between cycles; as noted in the Preliminaries, CAEM achieves that bound with a low-rank adapter on a frozen backbone, a strict training-pool filter and a retention guard rather than with an explicit penalty. The low-rank choice is reviewed next.

The choice to fine-tune through a low-rank adapter rather than the full backbone is a separate and equally load-bearing decision in the self-improvement loop. Hu et al.~\cite{hu2022lora} introduced low-rank adaptation as a parameter-efficient alternative to full fine-tuning that injects trainable low-rank matrices into the weight matrices of a frozen backbone; the original paper adapts only the attention projections and reports quality on par with full fine-tuning at less than one percent of the trainable parameter count on GPT-3 175B, and later work extends the adapters to the feed-forward projections. The choice between low-rank adaptation and full fine-tuning has empirical consequences for continual learning that the recent literature has measured directly. Biderman et al.~\cite{biderman2024lora} compare low-rank adaptation with full fine-tuning for Llama-2-7B on code and mathematics, sweeping ranks from $16$ to $256$, and report two findings relevant to CAEM's setting: low-rank adaptation learns less of the target domain than full fine-tuning, but it forgets less of the base model's general capability, measured as average HellaSwag, ARC-Challenge, and WinoGrande accuracy. In code instruction tuning on Magicoder data, for example, full fine-tuning falls to 0.414 against 0.509 for low-rank adaptation at rank sixty-four. Forgetting grows with training duration in both cases, and in mathematics low-rank adaptation at rank 256 forgets nearly as much as full fine-tuning, so higher rank trades retention for fit. The QLoRA work of Dettmers et al.~\cite{dettmers2023qlora} establishes the second decisive practical point: applying low-rank adaptation to all linear transformer-block layers, both attention and feed-forward, rather than to the attention projections alone is the most critical configuration choice for matching full fine-tuning, whereas the rank itself has little effect once all layers are adapted. CAEM's adapter target set spans every linear projection inside the transformer block (the four attention projections together with the three feed-forward projections of the gated multi-layer perceptron), which is the all-linear-layer target set that this literature converges on, fitted at rank thirty-two with $\alpha=2r$; with this choice the effective scale $\alpha/r$ stays at two, so the rank-driven gradient collapse that motivates the rank-stabilised alternative for fixed-$\alpha$ adapters does not arise, and rank thirty-two is far below the large ranks at which the theory of Kalajdzievski~\cite{kalajdzievski2023rslora} predicts instability for a constant scale. Low-rank adaptation is also what lets the self-improvement loop fit on a single consumer-grade graphics processor; full fine-tuning at the same backbone size would not fit. The within-adapter analogue of catastrophic forgetting under sequential fine-tuning, documented by Wang et al.~\cite{wang2023olora} and orthogonal-decomposition fixes for it, is registered as future work; the present thesis combines the bounded low-rank-adapter parameter envelope on a frozen backbone, the multi-modal retention guard with adapter rollback, and the strict training threshold above the storage cut-point as the three-mechanism defence against iterative-adapter drift across the realised cycle trajectory.

Natarajan et al. \cite{natarajan2013learning} provided theoretical foundations for learning with noisy labels at NeurIPS 2013. Their unbiased estimator method constructs proxy loss functions recovering optimal classifiers despite label corruption. The key insight is that if the noise process is known or can be estimated (e.g., what fraction of positive labels are actually negative), modified loss functions achieve consistent parameter estimation. On a synthetic non-separable dataset their methods achieve more than 88\% accuracy even when 40\% of the labels are corrupted, and they are competitive with other methods for learning under label noise. 

For CAEM, this framework explains why training on an imperfect pool can still help: no verifier is perfect, so some stored episodes carry wrong answers. Natarajan et al.'s guarantees, however, assume a loss corrected for the noise rates, which CAEM does not apply (\Cref{thm:bayes-purity}). CAEM instead lowers the noise rate directly, by admitting to the training pool only episodes above a stricter threshold $\tau_{\text{train}} > \tau_{\text{store}}$, so that the training pool is purer than the memory as a whole. The realised purity of the training pool is reported in \Cref{ch:evaluation}.

When Zelikman et al. \cite{NEURIPS2022_639a9a17} introduced STaR (Self Taught Reasoner) at NeurIPS 2022, they showed successful bootstrapping using iterative self improvement. STaR loop generates rationales for problems, filters for correct answers, fine tunes on correct rationale-answer pairs and then repeats. A key innovation is rationalisation. For problems the model fails to answer correctly, it is given the correct answer as a hint and generates a rationale that leads to it. This data augmentation adds training signal from problems the model could not yet solve, which the basic loop discards.

On CommonsenseQA, STaR improved over a few-shot baseline by 35.9 percentage points (72.5\% against 36.6\%), showing that bootstrapping from a model's own outputs works when coupled with appropriate filtering. The parallel to CAEM is direct: generate predictions, verify them, fine-tune on verified examples, and repeat. CAEM extends STaR in two ways. First, simple correctness filtering is replaced by the nine-signal verifier, whose calibrated composite $\ustored$ drives the four-outcome storage decision. Second, verified items persist in episodic memory for future retrieval rather than being consumed only as training data.

Recent work on intrinsic self-correction reveals important limitations. Huang et al.\ \cite{huang2024cannot} found at ICLR 2024 that prompting language models to review and correct their own reasoning does not improve accuracy without external feedback, and at times degrades it. Kamoi et al.\ \cite{kamoi-etal-2024-llms} analysed at TACL 2024 when language models can correct their mistakes and found that intrinsic self-correction often fails. The bottleneck is reliable feedback: a model that produces an incorrect first answer usually cannot generate a reliable correction signal without external validation.

This limitation motivates CAEM's design. Instead of asking the model to review itself, CAEM uses external verification: an entailment judge against retrieved evidence, agreement across resampled chains, and recalibrated internal signals. These supply the external feedback that the self-correction literature identifies as a precondition for useful self-improvement, and because their failure modes differ, a calibrated combination of them is expected to be more reliable than any one of them or than the model's own self-evaluation. \Cref{ch:evaluation} reports how well the verifier separates correct from incorrect answers.

The self-improvement literature also reports successes: when external feedback is available, for example code execution results, formal verification or retrieved sources, models can use it to correct themselves. CAEM's retrieval-augmented tier supplies feedback of this kind, since the retrieved evidence both grounds the generation and gives the verifier something to check the answer against.

Domain adaptation research provides additional relevant findings. Gururangan et al. \cite{gururangan-etal-2020-dont} demonstrated at ACL 2020 that continued pre-training on domain-specific unlabelled data (domain-adaptive pre-training) followed by task-specific fine-tuning outperforms direct fine-tuning on limited labelled data. While CAEM does not perform domain-adaptive pre-training due to computational constraints, the finding supports iterative refinement: progressive training on accumulated verified examples enables gradual domain adaptation.

The distinction between CAEM's approach and standard continual learning deserves emphasis. Traditional continual learning addresses learning sequential tasks A, B, C without forgetting previous tasks. CAEM's scenario differs: the task remains constant (question answering), but the model progressively improves at that task through accumulated verified examples. This setting is closer to curriculum learning and iterative refinement than task-sequential continual learning, though catastrophic forgetting concerns still apply since fine-tuning on domain-specific verified examples risks degrading general capabilities.

Several 2024--2025 contributions tighten the theoretical and architectural framing of self-improvement in ways directly relevant to CAEM's theoretical analysis and to its retention guard. Hosseini et al.~\cite{hosseini2024vstar} introduced V-STaR at COLM 2024, training a DPO verifier over correct and incorrect self-generated solutions and using it to rank test-time candidates; CAEM generalises V-STaR by using the verifier for training-data admission rather than only for test-time ranking, and by adding the episodic memory and the retention guard. Song et al.~\cite{song2024gap} formalise the \emph{generation--verification gap}, the improvement obtained by reweighting a model's generations with its own verification, and show empirically that self-improvement requires a positive gap, that small models fail to achieve one, and that iterative self-improvement saturates within two or three rounds; this is the conceptual motivation for CAEM's $\alpha > 1/2$ verifier-purity precondition in \Cref{sec:theory}. Yi et al.~\cite{fu2025collapse} prove theoretically, in linear regression, that a verifier filter over synthetic data prevents model collapse in self-consuming training loops and yields near-term improvement, while in the long run the model converges to the verifier's knowledge centre, so that gains plateau or reverse unless the verifier is unbiased; this is the strongest published safety anchor for CAEM's storage-gate-filtered SIL, and also its main caution. Das et al.~\cite{das2025retraining} prove at ICML 2025 that retraining a model on its own predicted hard labels increases population accuracy in a binary linearly-separable classification setting with randomly flipped labels, and show empirically that retraining only on samples where the prediction agrees with the given label helps further; this is the closest published analogue of CAEM's pool-purity inequality (\Cref{thm:purity}), and CAEM extends the idea from binary classification with self-predicted labels to open-ended generation with an external multi-signal verifier. Huang, Block, Foster et al.~\cite{huang2025sharpening} formalise self-improvement as a sharpening mechanism at ICLR 2025 and prove convergence-style amortisation-via-self-training bounds that are the closest published analogue of CAEM's bounded-band stationarity heuristic on the committed-cycle subsequence (\Cref{thm:convergence}). Lang, Sontag, and Vijayaraghavan~\cite{lang2024w2s} provide a complementary expansion-based bound at NeurIPS 2024 for when a strong student trained on weak-teacher pseudolabels outperforms the teacher; their bound requires expansion assumptions rather than a verifier-accuracy threshold and is therefore complementary to, rather than redundant with, CAEM's data-purity formulation.

These methods inform CAEM's self-improvement loop as follows. EWC motivates bounding drift between cycles, which CAEM achieves with a low-rank adapter on a frozen backbone, an approach the parameter-efficient fine-tuning literature shows forgets less than full fine-tuning. Noisy-label learning shows that learning from imperfect labels remains possible when the noise is bounded, which motivates the stricter training threshold. STaR supplies the generate, filter, fine-tune and repeat cycle. The self-correction literature motivates external verification rather than self-review. Finally, the retention guard, which rolls the adapter back but never the memory, turns forgetting into a detected and rejected event rather than a silent degradation. The trajectory actually achieved is reported in \Cref{ch:evaluation}.

\Cref{tab:positioning} positions CAEM against the closest published systems on the architectural axes that distinguish them: whether the system stores a per-entry confidence at insertion, what signal the router uses, what filter the training-data admission applies, what regulariser preserves general capability under self-improvement, and whether a formal precision or convergence guarantee is supplied. No single prior system shares CAEM on every axis simultaneously; the table makes the differentiation explicit.

\begin{table}[!htbp]
\centering
\footnotesize
\begin{tabular}{@{}p{2.4cm}p{1.7cm}p{2.6cm}p{2.5cm}p{2.0cm}p{2.0cm}@{}}
\toprule
\textbf{System} & \textbf{Stored conf.} & \textbf{Routing signal} & \textbf{Training filter} & \textbf{CL guard} & \textbf{Formal guarantee} \\
\midrule
\textbf{CAEM} (this work) & yes (calibrated $\ustored$) & combined sim $+$ $\ustored$, $\upre$ override & 9-signal cal-prob composite $+$ fixed-threshold gate (per-bench shrinkage; LoRA SIL) & bounded LoRA adapter envelope $+$ multi-modal retention rollback & calibration-consistent precision floor; purification condition; trajectory heuristics \\
A-MEM~\cite{xu2025amem} & no (textual tags) & similarity $+$ link evolution & --- & --- & --- \\
Larimar~\cite{larimar2024das} & no (key-value addr.) & key match & one-shot edit & --- & --- \\
Adaptive-RAG~\cite{jeong2024adaptive} & no & query complexity classifier & --- & --- & --- \\
V-STaR~\cite{hosseini2024vstar} & no & --- & DPO verifier (test-time rank) & --- & --- \\
Mohri-Hashimoto~\cite{mohri2024conformal} & no & --- & split-conformal sub-claim filter & --- & conformal factuality \\
MiniCheck~\cite{tang2024minicheck} & no & --- & NLI judge (fact verification) & --- & --- \\
\bottomrule
\end{tabular}
\caption[Positioning of CAEM against the closest prior systems.]{Positioning of CAEM against the closest prior systems. Insertion-time stored confidence, combined-score routing on a shared memory, the calibrated-probability composite at a fixed-threshold gate, the LoRA self-improvement primitive with an adapter-rollback guard against a multi-modal retention probe panel, and the calibration-consistent precision floor with its purification condition together constitute CAEM's distinguishing combination.}
\label{tab:positioning}
\end{table}

\subsection{External Baselines and Competitive Positioning}
\label{sec:external-baselines}

The related work reviewed above clusters into three families relevant to CAEM's evaluation: prompt-level reasoning methods, retrieval-augmented generation methods, and training-time self-improvement methods, to which a confidence-calibration method is added. \Cref{ch:evaluation} runs representatives of each as external baselines, so that every mechanism CAEM adds to a plain language model is compared with at least one prior system that uses that mechanism alone. This subsection states what each baseline is and what it isolates; the numerical results appear in \Cref{ch:evaluation}.

The prompt-level family contributes three baselines. The first, \textbf{zero-shot Qwen-2.5-3B-Instruct}, uses the same decoder-only backbone CAEM fine-tunes, with no CoT prefix, no retrieval, and no verification. This is the empirical floor: any gain CAEM reports must exceed the zero-shot score on the same backbone, otherwise the gain is attributable to model capacity rather than to the CAEM mechanism. The second, \textbf{Chain-of-Thought} prompting~\cite{wei2022chain}, adds the standard zero-shot ``Let's think step by step'' trigger~\cite{kojima2022large} without any further modification. CoT isolates the contribution of intermediate reasoning on the same backbone and controls for the possibility that CAEM's gains reflect better rationale generation rather than the verifier-memory architecture. The third, \textbf{Five-shot Chain-of-Thought}, prepends five in-context $(\text{question}, \text{reasoning}, \text{answer})$ demonstrations from a benchmark training split before the CoT trigger and is the strongest prompting-only comparison against CAEM's retrieval-free Tier-2.

The retrieval family contributes three baselines. \textbf{DPR-RAG}~\cite{karpukhin-etal-2020-dense, NEURIPS2020_6b493230} retrieves the top $k=5$ Wikipedia passages from a dense DPR index, concatenates them into a 384-token context, and generates from the same Qwen-2.5-3B-Instruct backbone through the Tier-3 prompt builder; this isolates the contribution of external grounding without any verification or memory. \textbf{CoT$+$DPR-RAG} combines the CoT trigger with DPR retrieval and isolates the contribution of reasoning over retrieved evidence. \textbf{FLARE}~\cite{jiang-etal-2023-active} extends DPR-RAG with active retrieval: rather than retrieving once at the start of generation, FLARE generates one sentence at a time and triggers a retrieval whenever the token-level confidence on the upcoming sentence falls below a registered threshold, with retrieved passages prepended to the generation context for the next sentence; this isolates the contribution of selective and step-conditioned retrieval over single-pass retrieval. A fourth retrieval system, \textbf{Self-RAG}~\cite{asai2024selfrag}, is acknowledged as the most ambitious retrieve-and-critique pipeline in the literature but is included only as a citation-level reference rather than as a replicated numerical baseline, because Self-RAG's critic tokens are tied to a LLaMA-family backbone and retargeting them to Qwen-2.5-3B-Instruct would change enough of the system to invalidate the architectural comparison.

The training-time family contributes one baseline. \textbf{Vanilla fine-tuning} fine-tunes all weights of Qwen-2.5-3B-Instruct on raw question-answer pairs from the three training benchmarks for the same number of cycles as CAEM, with no verifier filter, no retention guard and no episodic memory. It isolates the contribution of CAEM's curation and retention machinery: if vanilla fine-tuning degrades on the retention probes or across cycles while CAEM does not, that machinery is responsible. The confidence-calibration family contributes one baseline. \textbf{Semantic entropy}~\cite{farquhar2024semantic} samples several generations at a raised temperature, clusters them into groups of equivalent meaning, and uses the entropy of the cluster distribution as the per-query uncertainty; the baseline runs the same backbone as the zero-shot system and abstains when the entropy is high, which makes it the most direct external comparator to CAEM's calibrated composite. \textbf{STaR}~\cite{NEURIPS2022_639a9a17}, discussed above as the closest prior method in spirit, is not run: its rationalisation step and binary correctness filter differ enough from CAEM's verifier and four-outcome decision that a comparison would mix architectural contribution with rationalisation quality.

Taken together, the eight external baselines (B1 to B8 in \Cref{sec:comp-baselines}) span the prior systems that share at least one mechanism with CAEM: three isolate the prompt-level component (B1 zero-shot, B2 CoT, B5 five-shot CoT), three the retrieval component (B3 DPR-RAG, B4 CoT$+$DPR-RAG, B7 FLARE), one the training-time component (B6 vanilla fine-tuning), and one a sampling-based uncertainty estimate on the same backbone (B8 semantic entropy). Neither Self-RAG nor STaR is run, so neither enters the significance tests. No single baseline combines all four families; CAEM's contribution is their combination, which the comparison in \Cref{ch:evaluation} tests against each family in isolation.

\subsection{Evaluation Benchmarks}

Rigorous evaluation requires benchmarks testing diverse reasoning capabilities with objective correctness criteria. This subsection reviews the five primary factual-QA benchmarks that constitute CAEM's locked evaluation panel (FEVER, TriviaQA, CommonsenseQA training, and TruthfulQA, StrategyQA transfer) and the MMLU general-knowledge probe that sits inside the multi-modal retention guard, explaining how each tests a different aspect of hallucination reduction. Natural Questions is also described for completeness: it and HotpotQA were evaluated at cycle zero and dropped before the main run because the base generator answered too few of their questions correctly for the storage gate to build a usable training pool (\Cref{sec:scope-research}).

Thorne et al.~\cite{thorne-etal-2018-fever} created the FEVER dataset at NAACL 2018 for fact verification, providing 185{,}445 claims generated from Wikipedia requiring classification as SUPPORTS, REFUTES, or NOT ENOUGH INFO relative to Wikipedia evidence. Claims are designed to test factual knowledge; an illustrative pair (not taken from the dataset) is ``Barack Obama was the 44th President of the United States'' (SUPPORTS) and ``The Beatles formed in Manchester'' (REFUTES, since they formed in Liverpool). The NOT ENOUGH INFO category is particularly valuable, testing whether systems recognise insufficient evidence rather than hallucinating conclusions. FEVER directly evaluates the external-grounding family of CAEM's verifier, because the task structure matches the premise-hypothesis entailment check that produces $p_{\text{ground}}$ and $p_{\text{entail}}$.

Joshi et al.~\cite{joshi-etal-2017-triviaqa} introduced TriviaQA at ACL 2017, providing about 95{,}000 trivia question-answer pairs authored by trivia enthusiasts, which with on average six supporting documents each from the web and Wikipedia give over 650{,}000 question-answer-evidence triples. Questions span an unusually wide topical range (history, science, geography, entertainment) and the provided evidence is noisy (only a subset of the retrieved documents actually answers the question), making TriviaQA a realistic test of retrieval-augmented generation under distractor documents. For CAEM, TriviaQA tests whether the Tier-3 retrieval-augmented path can ground answers against noisy evidence and whether the mean-match grounding signal $p_{\text{ground\_mean}}$ remains calibrated when retrieval quality varies across queries.

Kwiatkowski et al.~\cite{kwiatkowski-etal-2019-natural} introduced Natural Questions (NQ) at TACL 2019, providing real anonymised queries to the Google search engine paired with Wikipedia-page long answers and span-level short answers. Unlike synthetic benchmarks, NQ captures the distribution of questions users actually ask, with the accompanying variability in specificity, ambiguity, and presupposition. NQ is the dominant open-domain QA benchmark in recent retrieval-augmented-generation literature~\cite{NEURIPS2020_6b493230, karpukhin-etal-2020-dense}. CAEM evaluated NQ at cycle zero and dropped it before the main run: the base generator answered only $17$ percent of its questions correctly, too few for the storage gate to build a usable training pool. Its training split is reused only to train the retrieval-utility classifier.

Lin et al.~\cite{lin-etal-2022-truthfulqa} designed TruthfulQA specifically to measure how models mimic human falsehoods, providing 817 questions across 38 categories in health, law, finance, politics, and conspiracies. Questions are adversarially constructed to elicit common misconceptions: ``What happens if you crack your knuckles?'' where the false but common answer is ``you get arthritis'' and the true answer is ``nothing harmful happens; no proven link to arthritis.'' The inverse scaling finding (larger models often less truthful) makes TruthfulQA the most challenging benchmark for parametric scaling approaches. Human performance reaches 94\% truthful and 87\% truthful-and-informative, while the best GPT-3 175B configuration achieves 58\% and 21\% respectively. For CAEM, TruthfulQA was registered as the benchmark where the sample-set family of signals was expected to pay off, on the reasoning that a misconception would make the model alternate between the correct answer and the misconception across stochastic decodes, which the sample-set agreement signals $s_{\text{avg}}$ and $p_{\text{entail}}$ would detect. The literature predicts otherwise: semantic entropy targets arbitrary, sample-to-sample errors and explicitly does not cover models that are consistently wrong because of common misconceptions, as Farquhar et al.~\cite{farquhar2024semantic} state, since such a model gives the same wrong answer on every decode and agreement-based signals then read it as confident. The realised results agree with the literature rather than with the registered expectation: entries stored on TruthfulQA at cycle zero are correct only 52.2\% of the time, against 75 to 82\% on the training panel (\Cref{tab:cross-bench-summary}), and the semantic-entropy baseline is among the highest-scoring baselines on TruthfulQA exact match rather than being penalised for heavy abstention (\Cref{ch:evaluation}). TruthfulQA is therefore the benchmark that exposes the blind spot of agreement-based verification.

Geva et al.~\cite{10.1162/tacl_a_00370} introduced StrategyQA at TACL 2021, presenting 2{,}780 questions requiring implicit reasoning where decomposition steps are not mentioned in questions. For example, the dataset question ``Could a llama birth twice during the War in Vietnam?'' requires knowing Vietnam War duration (implicit: 1955--1975) and llama gestation period (implicit: approximately 11 months), then computing whether two births fit the timeframe. Human accuracy reaches 87\% while best models achieve approximately 66\%. The implicit reasoning requirement tests whether models genuinely understand underlying concepts or merely perform shallow pattern matching. For CAEM, StrategyQA evaluates the sample-set agreement signal $s_{\text{avg}}$ because implicit reasoning should converge across multiple chains if the underlying world-knowledge is consistent; divergence across chains is a strong hallucination signal specifically on this benchmark.

Hendrycks et al.~\cite{hendrycks2021measuring} introduced MMLU in 2021, a 57-subject multiple-choice benchmark covering humanities, STEM, social sciences, and professional domains. MMLU is used in CAEM as one of three components of the \textit{multi-modal retention probe panel}, not as a target benchmark. The fine-tune is not optimised for MMLU; instead, MMLU is probed before the first cycle (the pristine baseline) and after every cycle's fine-tune, alongside two further probes drawn from held-out test slices of the training panel: an open-text factoid probe on the held-out TriviaQA test fold, and a commonsense multiple-choice probe on the held-out CommonsenseQA test fold. The cycle aborts and rolls the adapter back to the previous snapshot whenever the worst per-probe ratio against the pristine baseline falls below the fixed floor $\rho_{\min} = 0.93$; the worst-probe rule replaces an earlier single-probe rule because a general-knowledge multiple-choice probe alone missed the failure mode where MCQ accuracy stayed high while open-text generation degraded silently. The episodic memory is preserved across rollbacks. The choice of MMLU as one of the three probes is deliberate: catastrophic forgetting manifests most visibly as erosion of capability on tasks the model was not trained on this cycle, so a broad out-of-distribution probe is a more sensitive trigger than in-distribution QA, and the held-out training-bench test folds add the asymmetric-drift sensitivity that the broad probe alone would miss.

The benchmarks are complementary. FEVER tests entailment-based grounding, including the NOT ENOUGH INFO case where the right behaviour is to recognise missing evidence. TriviaQA tests retrieval-augmented generation under noisy evidence. CommonsenseQA tests commonsense reasoning over short multiple-choice options. TruthfulQA tests agreement-based and relevance signals on imitative falsehoods, as a transfer-only benchmark held out from the SIL training pool. StrategyQA tests sample-set agreement on implicit multi-hop reasoning, also as a transfer-only benchmark. MMLU, the open-text TriviaQA test fold and the multiple-choice CommonsenseQA test fold close the catastrophic-forgetting loop through the retention probe panel. Together they cover a range of factual reasoning tasks, so that results on the panel show whether CAEM's effects carry across reasoning types rather than fitting one task.

\section{Summary of Key Findings}

The literature reviewed in this chapter supports five findings that shape CAEM.

First, hallucination is persistent and is not solved by scale. Rates vary widely by domain, up to 91.4\% of chatbot-generated references in one systematic-review study \cite{info:doi/10.2196/53164}; on TruthfulQA the best GPT-3 configuration was truthful and informative on only 21\% of questions, with larger models often less truthful \cite{lin-etal-2022-truthfulqa}; and FActScore finds only 58\% of ChatGPT's atomic facts supported \cite{min-etal-2023-factscore}. This motivates changing the structure around the model rather than relying on scale.

Second, no single verification or confidence signal is reliable across tasks. Verbalised confidence discriminates weakly \cite{xiong2023can}, modern networks are systematically overconfident \cite{pmlr-v70-guo17a}, preference-tuned models are miscalibrated until rescaled \cite{kadavath2022language, tian2023just}, and uncertainty methods that do well on open generation do poorly on constrained outputs and vice versa \cite{vashurin2024polygraph}. Combining calibrated signals has published precedent in the logistic multi-scoring of \cite{valentin2024hallucination}, and generic NLI models transfer poorly to language-model hallucinations \cite{sinha2025semantic}, which favours a judge trained on claim-support data \cite{tang2024minicheck}. CAEM therefore combines nine recalibrated signals into one composite.

Third, external memory and retrieval improve factual accuracy and can be updated without retraining \cite{graves2016differentiable, NEURIPS2020_6b493230}, but standard retrieval-augmented generation neither verifies what it retrieves nor accumulates verified answers, and retrieval does not help every query \cite{mallen2023whennot, wang2023skr}. CAEM therefore stores only verified answers, routes on similarity combined with stored confidence, and decides per query whether retrieval is worth using.

Fourth, self-improvement works when it receives reliable external feedback and fails when a model only reviews itself \cite{NEURIPS2022_639a9a17, huang2024cannot, kamoi-etal-2024-llms}; verifier-filtered retraining can help but saturates within a few rounds \cite{song2024gap, fu2025collapse}. CAEM therefore filters its training data through an external verifier, with a stricter threshold than its memory uses.

Fifth, fine-tuning causes forgetting, which low-rank adaptation reduces but does not remove \cite{hu2022lora, biderman2024lora}, and weight-regularisation penalties add little once the backbone is pre-trained \cite{mehta2023empirical, wu2022pretrained}. CAEM therefore relies on a low-rank adapter on a frozen backbone and on a retention guard that rolls back the adapter, but never the memory, when held-out accuracy falls.

The five-benchmark panel and the three-probe retention panel exercise these mechanisms on different failure modes, and the baselines of \Cref{sec:external-baselines} each isolate one family of prior methods. No prior system combines all of these elements (\Cref{tab:positioning}), which is the gap CAEM addresses.

\section{Chapter signpost}
\label{sec:ch2-signpost}

The preliminaries section and the literature review have established the technical foundations on which the architecture rests and have located CAEM's joint intervention against the verification, retrieval, memory-augmented, and continual-learning prior work whose individual contributions do not close the deployment-safety gap registered in \Cref{ch:introduction}. \Cref{ch:requirements} translates the architectural opportunity surfaced by the literature into a registered specification: seven functional requirements paired with named architectural components, five non-functional requirements with empirical anchors, six production-parity design commitments, and the societal, environmental, ethical, standards, project-management, risk, and economic readings that bound the engineering envelope. The specification is then discharged end-to-end by the methodology chapter that follows.